% GENERATED FILE. Edit canonical lessons, then run npm run generate:books.
% canonical-source-hash: fnv1a64:6a5c3fb501e6413c
% canonical-lessons: GU-C399-vilamb, GU-C399-sambhlavu, GU-C399-elan, GU-C399-talvu, GU-C399-dokiyu

\chapter{Delays and Announcements}
\label{ch:gu-delays-and-announcements}
\hlchaptermodality{399}

\begin{chapteropening}
\textbf{By the end of this chapter:} \emph{I can understand an announcement about a late train.}

Complete the last lesson of chapter 399: I can understand an announcement about a late train.
\end{chapteropening}

\section[\texorpdfstring{\gu{વિલંબ} (vilamb)}{વિલંબ (vilamb)}]{\gu{વિલંબ} (vilamb) — a delay (formal)}

\label{lesson:GU-C399-vilamb}



\begin{quote}
Before the new one: say the Gujarati for a jolt, then the Gujarati for to rattle.
\end{quote}

\subsection*{You'll want to know: \gu{વિલંબ}}
\textbf{\gu{વિલંબ}} — \emph{vilamb} — \textquotedblleft{}a delay (formal)\textquotedblright{}.

\begin{quote}
\gu{ધુમ્મસને} \gu{લીધે} \gu{ટ્રેનમાં} \gu{બે} \gu{કલાકનો} \gu{વિલંબ} \gu{છે}.
\end{quote}

\begin{quote}
\emph{dhummasne līdhe ṭrenmā̃ be kalākno vilamb che.}
\end{quote}

\begin{quote}
Because of the fog the train is two hours late.
\end{quote}

A masculine noun.

\begin{cousinweb}
\textbf{\gu{વિલંબ}} (\emph{vilamb}) is a tatsama, Sanskrit \emph{vilamba}, \textquotedblleft{}hanging down\textquotedblright{}, then \textquotedblleft{}slowness, delay\textquotedblright{} (Monier-Williams); station announcements use it where speech says \gu{મોડું}.
\end{cousinweb}

\begin{grammarlens}[title={What you've built}]
One more word for getting around.
\end{grammarlens}

\subsection*{Your turn}
\begin{itemize}
\raggedright
  \item \emph{Say it:} \emph{vilamb}
  \item \emph{Say it:} \emph{vilamb}, once more
  \item \emph{Say it:} \emph{dhummasne līdhe ṭrenmā̃ be kalākno vilamb che.} — Because of the fog the train is two hours late.
  \item \emph{Recall:} say the Gujarati for a jolt, then the Gujarati for to rattle, then say \emph{vilamb} again
\end{itemize}

\begin{tcolorbox}[breakable,colback=teal!4,colframe=teal!35!black,title={Before you move on}]
What does \gu{વિલંબ} mean? (\textquotedblleft{}A delay (formal)\textquotedblright{}.) Say it once more.
\end{tcolorbox}
\section[\texorpdfstring{\gu{સંભળાવું} (sambhḷāvũ)}{સંભળાવું (sambhḷāvũ)}]{\gu{સંભળાવું} (sambhḷāvũ) — to be heard; to be audible}

\label{lesson:GU-C399-sambhlavu}



\begin{quote}
Before the new one: say the Gujarati for to rattle, then the Gujarati for a delay.
\end{quote}

\subsection*{You'll want to know: \gu{સંભળાવું}}
\textbf{\gu{સંભળાવું}} — \emph{sambhḷāvũ} — \textquotedblleft{}to be heard; to be audible\textquotedblright{}.

\begin{quote}
\gu{દૂરથી} \gu{ટ્રેનની} \gu{સીટી} \gu{સંભળાય} \gu{છે}.
\end{quote}

\begin{quote}
\emph{dūrthī ṭrennī sīṭī sambhḷāy che.}
\end{quote}

\begin{quote}
The train's whistle can be heard from far away.
\end{quote}

Past (it …): \gu{તે} \gu{અવાજ} \gu{સંભળાયો}, \emph{te avāj sambhḷāyo}.

\begin{cousinweb}
\textbf{\gu{સંભળાવું}} (\emph{sambhḷāvũ}) is the intransitive partner of \gu{સાંભળવું}, \textquotedblleft{}to listen, to hear\textquotedblright{}, which goes back to \emph{*saṁbhalati}, \textquotedblleft{}attends to\textquotedblright{}, a form CDIAL reconstructs (CDIAL 12962): what you hear, \gu{સાંભળવું}; what can be heard, \gu{સંભળાવું}.
\end{cousinweb}

\begin{grammarlens}[title={What you've built}]
One more word for getting around.
\end{grammarlens}

\subsection*{Your turn}
\begin{itemize}
\raggedright
  \item \emph{Say it:} \emph{sambhḷāvũ}
  \item \emph{Say it:} \emph{sambhḷāvũ}, once more
  \item \emph{Say it:} \emph{dūrthī ṭrennī sīṭī sambhḷāy che.} — The train's whistle can be heard from far away.
  \item \emph{Recall:} say the Gujarati for to rattle, then the Gujarati for a delay, then say \emph{sambhḷāvũ} again
\end{itemize}

\begin{tcolorbox}[breakable,colback=teal!4,colframe=teal!35!black,title={Before you move on}]
What does \gu{સંભળાવું} mean? (\textquotedblleft{}To be heard; to be audible\textquotedblright{}.) Say it once more.
\end{tcolorbox}
\section[\texorpdfstring{\gu{એલાન} (elān)}{એલાન (elān)}]{\gu{એલાન} (elān) — a public announcement}

\label{lesson:GU-C399-elan}



\begin{quote}
Before the new one: say the Gujarati for a delay, then the Gujarati for to be heard.
\end{quote}

\subsection*{You'll want to know: \gu{એલાન}}
\textbf{\gu{એલાન}} — \emph{elān} — \textquotedblleft{}a public announcement\textquotedblright{}.

\begin{quote}
\gu{સ્ટેશન} \gu{પર} \gu{એલાન} \gu{સંભળાયું} \gu{કે} \gu{ટ્રેન} \gu{મોડી} \gu{છે}.
\end{quote}

\begin{quote}
\emph{sṭeśan par elān sambhḷāyũ ke ṭren moḍī che.}
\end{quote}

\begin{quote}
An announcement at the station said the train was late.
\end{quote}

A neuter noun.

\begin{cousinweb}
\textbf{\gu{એલાન}} (\emph{elān}) is usually traced to Arabic \emph{iʿlān}, \textquotedblleft{}a declaration, a making known\textquotedblright{}, which reached Gujarati through Persian.
\end{cousinweb}

\begin{grammarlens}[title={What you've built}]
One more word for getting around.
\end{grammarlens}

\subsection*{Your turn}
\begin{itemize}
\raggedright
  \item \emph{Say it:} \emph{elān}
  \item \emph{Say it:} \emph{elān}, once more
  \item \emph{Say it:} \emph{sṭeśan par elān sambhḷāyũ ke ṭren moḍī che.} — An announcement at the station said the train was late.
  \item \emph{Recall:} say the Gujarati for a delay, then the Gujarati for to be heard, then say \emph{elān} again
\end{itemize}

\begin{tcolorbox}[breakable,colback=teal!4,colframe=teal!35!black,title={Before you move on}]
What does \gu{એલાન} mean? (\textquotedblleft{}A public announcement\textquotedblright{}.) Say it once more.
\end{tcolorbox}
\section[\texorpdfstring{\gu{ટળવું} (ṭaḷvũ)}{ટળવું (ṭaḷvũ)}]{\gu{ટળવું} (ṭaḷvũ) — to be averted; to be put off}

\label{lesson:GU-C399-talvu}



\begin{quote}
Before the new one: say the Gujarati for to be heard, then the Gujarati for a public announcement.
\end{quote}

\subsection*{You'll want to know: \gu{ટળવું}}
\textbf{\gu{ટળવું}} — \emph{ṭaḷvũ} — \textquotedblleft{}to be averted; to be put off\textquotedblright{}.

\begin{quote}
\gu{પવનની} \gu{દિશા} \gu{ફરી} \gu{ગઈ}, \gu{એટલે} \gu{મોટું} \gu{તોફાન} \gu{ટળી} \gu{ગયું}.
\end{quote}

\begin{quote}
\emph{pavannī diśā pharī gaī, eṭle moṭũ tophān ṭaḷī gayũ.}
\end{quote}

\begin{quote}
The wind changed direction, so the big storm was averted.
\end{quote}

Past (it …): \gu{તે} \gu{તોફાન} \gu{ટળી} \gu{ગયું}, \emph{te tophān ṭaḷī gayũ}.

\begin{cousinweb}
\textbf{\gu{ટળવું}} (\emph{ṭaḷvũ}) is a tadbhava of Sanskrit \emph{ṭálati}, \textquotedblleft{}moves aside\textquotedblright{} (CDIAL 5450); \gu{ટાળવું}, to avoid, is its causative partner: what you put off, \gu{ટાળવું}; what passes off by itself, \gu{ટળવું}.
\end{cousinweb}

\begin{grammarlens}[title={What you've built}]
One more word for getting around.
\end{grammarlens}

\subsection*{Your turn}
\begin{itemize}
\raggedright
  \item \emph{Say it:} \emph{ṭaḷvũ}
  \item \emph{Say it:} \emph{ṭaḷvũ}, once more
  \item \emph{Say it:} \emph{pavannī diśā pharī gaī, eṭle moṭũ tophān ṭaḷī gayũ.} — The wind changed direction, so the big storm was averted.
  \item \emph{Recall:} say the Gujarati for to be heard, then the Gujarati for a public announcement, then say \emph{ṭaḷvũ} again
\end{itemize}

\begin{tcolorbox}[breakable,colback=teal!4,colframe=teal!35!black,title={Before you move on}]
What does \gu{ટળવું} mean? (\textquotedblleft{}To be averted; to be put off\textquotedblright{}.) Say it once more.
\end{tcolorbox}
\section[\texorpdfstring{\gu{ડોકિયું} (ḍokiyũ)}{ડોકિયું (ḍokiyũ)}]{\gu{ડોકિયું} (ḍokiyũ) — a peep; a quick look}

\label{lesson:GU-C399-dokiyu}



\begin{quote}
Before the new one: say the Gujarati for a public announcement, then the Gujarati for to be averted.
\end{quote}

\subsection*{You'll want to know: \gu{ડોકિયું}}
\textbf{\gu{ડોકિયું}} — \emph{ḍokiyũ} — \textquotedblleft{}a peep; a quick look\textquotedblright{}.

\begin{quote}
\gu{મેં} \gu{બારીમાંથી} \gu{બહાર} \gu{ડોકિયું} \gu{કર્યું}.
\end{quote}

\begin{quote}
\emph{mẽ bārīmā̃thī bahār ḍokiyũ karyũ.}
\end{quote}

\begin{quote}
I took a peep out of the window.
\end{quote}

A neuter noun.

\begin{cousinweb}
\textbf{\gu{ડોકિયું}} (\emph{ḍokiyũ}) is built on \gu{ડોક}, the neck, which CDIAL lists under \emph{*ḍōka}, \textquotedblleft{}head\textquotedblright{} (CDIAL 5566): you stretch your neck out to look.
\end{cousinweb}

\begin{grammarlens}[title={What you've built}]
One more word for getting around.
\end{grammarlens}

\subsection*{Your turn}
\begin{itemize}
\raggedright
  \item \emph{Say it:} \emph{ḍokiyũ}
  \item \emph{Say it:} \emph{ḍokiyũ}, once more
  \item \emph{Say it:} \emph{mẽ bārīmā̃thī bahār ḍokiyũ karyũ.} — I took a peep out of the window.
  \item \emph{Recall:} say the Gujarati for a public announcement, then the Gujarati for to be averted, then say \emph{ḍokiyũ} again
\end{itemize}

\begin{tcolorbox}[breakable,colback=teal!4,colframe=teal!35!black,title={Before you move on}]
What does \gu{ડોકિયું} mean? (\textquotedblleft{}A peep; a quick look\textquotedblright{}.) Say it once more.
\end{tcolorbox}
